\section*{Hoofdstuk \ref{ch:b3:fouriertransform} --- De
fouriertransformatie}

\begin{solution}{exo:b3:fouriertransform:1}
$\widehat{\mathbf 1_{\intcc{-a}a}}(\xi) =
\int_{-a}^a\eu^{-\iu\xi x}\dd x = \frac{2\sin(a\xi)}{\xi}$
(waarde $2a$ in $0$).
$\widehat{\eu^{-a\abs x}}(\xi) = \int_0^\infty\eu^{-(a +
\iu\xi)x} + \eu^{-(a - \iu\xi)x}\,\dd x = \frac1{a + \iu\xi} +
\frac1{a - \iu\xi} = \frac{2a}{a^2 + \xi^2}$.
Tent: $\max(0, 1 - \abs x) = \mathbf 1_{\intcc{-1/2}{1/2}} *
\mathbf 1_{\intcc{-1/2}{1/2}}$, dus is haar getransformeerde
$\bigl(\frac{2\sin(\xi/2)}\xi\bigr)^2 =
\bigl(\frac{\sin(\xi/2)}{\xi/2}\bigr)^2$.
Laatste: $\frac{2a}{a^2+\xi^2} \in L^1$, dus geeft de inversie
(\cref{thm:b3:fouriertransform:inversion}(2)) toegepast op
$\eu^{-a\abs x}$, na hernoeming van de veranderlijken,
\[
\widehat{\Bigl(\frac1{x^2 + a^2}\Bigr)}(\xi) =
\frac{\pi}{a}\,\eu^{-a\abs\xi} .
\]
\end{solution}

\begin{solution}{exo:b3:fouriertransform:2}
Uit \cref{prop:b3:fouriertransform:basic}:
$\widehat{f(\cdot - a)} = \eu^{-\iu a\xi}\hat f(\xi)$;
$\widehat{f\cos(b\cdot)} = \frac12\bigl(\hat f(\xi - b) + \hat
f(\xi + b)\bigr)$;
$\widehat{f(a\cdot + b)}(\xi) = \frac1a\,\eu^{\iu
b\xi/a}\,\hat f(\xi/a)$ ($a > 0$);
$\widehat{\overline{f(-\cdot)}} = \overline{\hat f}$;
$\widehat{f * f} = \hat f^2$. Op de gaussische functie
($\widehat{\eu^{-x^2}} = \sqrt\pi\eu^{-\xi^2/4}$) is elke regel
een controle van één regel --- bijvoorbeeld heeft
$\eu^{-(x-a)^2}$ als getransformeerde $\sqrt\pi\,\eu^{-\iu
a\xi}\eu^{-\xi^2/4}$, wat de rechtstreekse berekening (kwadraat
afsplitsen) bevestigt.
\end{solution}

\begin{solution}{exo:b3:fouriertransform:3}
(a) $h = \mathbf 1_{\intcc{-1}1}*\mathbf 1_{\intcc{-1}1}$ heeft
$\hat h = \bigl(\frac{2\sin\xi}\xi\bigr)^2 \in L^1$; inversie in
$x = 0$, waar $h(0) = \lambda(\intcc{-1}1\cap\intcc{-1}1) = 2$:
\[
2 = \frac1{2\pi}\int_\R\Bigl(\frac{2\sin\xi}\xi\Bigr)^2
\dd\xi
\ \Longrightarrow\
\int_\R\Bigl(\frac{\sin\xi}\xi\Bigr)^2\dd\xi = \pi ,
\]
in overeenstemming met $\int_0^\infty\frac{\sin^2}{\xi^2} =
\frac\pi2$ (\cref{pb:b3:lebesgue:1}).

(b) Plancherel voor $f = \eu^{-\abs x}$: $\int\abs{\hat f}^2 =
2\pi\int\abs f^2$ luidt $\int\frac{4\,\dd\xi}{(1 + \xi^2)^2} =
2\pi\int\eu^{-2\abs x}\dd x = 2\pi$: dus
$\int_\R\frac{\dd\xi}{(1+\xi^2)^2} = \frac\pi2$.
\end{solution}

\begin{solution}{exo:b3:fouriertransform:4}
(a) $g_t(x) = \frac1{2\sqrt{\pi t}}\eu^{-x^2/4t}$: volgens
\cref{ex:b3:fouriertransform:gaussian} met $a = \frac1{4t}$ is
$\hat g_t(\xi) = \frac1{2\sqrt{\pi t}}\sqrt{4\pi
t}\,\eu^{-t\xi^2} = \eu^{-t\xi^2}$.
(b) $\widehat{g_t * g_s} = \eu^{-t\xi^2}\eu^{-s\xi^2} =
\widehat{g_{t+s}}$, en de transformatie is injectief op $L^1$
(\cref{thm:b3:fouriertransform:inversion}(3)): dus $g_t * g_s =
g_{t+s}$.
(c) $\norm{g_t}_1 = 1$ (gaussische integraal);
$\norm{g_t}_2^2 = \frac1{4\pi t}\int\eu^{-x^2/2t}\dd x =
\frac{\sqrt{2\pi t}}{4\pi t} = \frac1{\sqrt{8\pi t}}$.
\end{solution}

\begin{solution}{exo:b3:fouriertransform:5}
Voor reële even $f$ is $\hat f(\xi) = \int f\cos(\xi x)\dd x$
(het sinusdeel valt weg): reëel en even. Oneven: $\hat f(\xi) =
-\iu\int f\sin(\xi x)$: oneven en zuiver imaginair. En $\hat f(0)
= \int f$: de totale massa. Is $f \geq 0$, dan is $\abs{\hat
f(\xi)} \leq \int\abs f = \int f = \hat f(0)$, zodat het supremum
in $0$ wordt bereikt. Voor een kansdichtheid is $\hat f(-\xi)$ de
karakteristieke functie van \cref{ch:b3:clt}: overal modulus
$\leq 1$, en $= 1$ in de oorsprong.
\end{solution}

\begin{solution}{exo:b3:fouriertransform:6}
(i) De injectiviteit is
\cref{thm:b3:fouriertransform:inversion}(3); de \omterm{def:b3:topology:continuity}{continuïteit} is
$\norm{\hat f}_\infty \leq \norm f_1$ (met waarden in $\mathcal
C_0$ volgens Riemann--Lebesgue); en $\mathcal C_0$ is gesloten in
de supremumnorm (uniforme limieten van functies die in oneindig
verdwijnen, verdwijnen in oneindig): dus banach.

(ii) Een \omterm{def:b3:topology:continuity}{continue} bijectie tussen banachruimten heeft een
\omterm{def:b3:topology:continuity}{continue} inverse (\cref{thm:b3:banach:openmapping}): er zou een
$C$ zijn met $\norm f_1 \leq C\norm{\hat f}_\infty$.

(iii) Zij $f_n(x) = \frac{\sin x}x\cdot\frac{\sin(x/n)}{x/n}$:
een product van twee $L^2$-functies, en $O(x^{-2})$ op oneindig,
dus $f_n \in L^1\cap L^2$. Omdat
$\bigl(\frac{\sin(ax)}{ax}\bigr)$ als $L^2$-getransformeerde
$\frac\pi a\mathbf 1_{\intcc{-a}a}$ heeft, geeft de
productformule $\widehat{gh} = \frac1{2\pi}\hat g * \hat h$
(geldig voor $g, h \in L^2$ met $gh \in L^1$; controleer haar op
schwartzfuncties met Fubini en breid uit met de $L^2$-continuïteit
van beide leden via Plancherel)
\[
\hat f_n = \frac1{2\pi}\,\bigl(\pi\mathbf
1_{\intcc{-1}1}\bigr) * \bigl(\pi n\,\mathbf
1_{\intcc{-1/n}{1/n}}\bigr):
\]
een trapezium van hoogte $\frac{\pi n}2\cdot\frac2n = \pi$: dus
$\norm{\hat f_n}_\infty = \pi$ voor elke $n$. Maar op $[1, n]$ is
$\frac{\sin(x/n)}{x/n} \geq \sin 1 > 0$, dus
\[
\norm{f_n}_1 \geq \sin 1\int_1^n\frac{\abs{\sin x}}x\dd x
\geq c\ln n
\]
(bogen tellen, als in \cref{thm:b3:banach:fourierdiverge}). De
afschatting $\norm{f_n}_1 \leq C\pi$ faalt voor grote $n$: dus
niet surjectief. (Het beeld is --- met argumenten van het type
Stone--Weierstrass --- een \omterm{def:b3:topology:interior}{dichte}, maar echte deelruimte van
$\mathcal C_0$.)
\end{solution}

\begin{solution}{exo:b3:fouriertransform:7}
Is $\hat f(\xi) = O(\abs\xi^{-k-1-\delta})$, dan is $\xi^j\hat f
\in L^1$ voor $0 \leq j \leq k$ (\omterm{def:b3:lebesgue:l1}{integreerbaar} op oneindig wegens
het verval, lokaal wegens de \omterm{def:b3:topology:continuity}{continuïteit} van $\hat f$). De
inversie (\cref{thm:b3:fouriertransform:inversion}(2))
vertegenwoordigt $f$ b.o.\ door $x \mapsto \frac1{2\pi}\int\hat
f(\xi)\eu^{\iu x\xi}\dd\xi$, en differentiëren onder de integraal
(dominanten $\abs{\xi^j\hat f}$) maakt die vertegenwoordiger
$\mathcal C^k$. Omgekeerd, voor $f \in \mathcal C_c^k$: iteratie
van \cref{prop:b3:fouriertransform:basic}(4) geeft
$(\iu\xi)^k\hat f = \widehat{f^{(k)}}$, dus $\abs{\hat f} \leq
\norm{f^{(k)}}_1\abs\xi^{-k}$. Tentfunctie: \omterm{def:b3:topology:continuity}{continu} met \omterm{def:b3:topology:compact}{compacte}
drager ($k = 0$: begrensde getransformeerde), en haar
getransformeerde $\sim \xi^{-2} = O(\abs\xi^{-0-1-1})$ levert
volgens de eerste richting een $\mathcal
C^0$-vertegenwoordiger --- allebei scherp: de tent is niet
$\mathcal C^1$, en haar getransformeerde vervalt niet sneller dan
$\xi^{-2}$.
\end{solution}

\begin{solution}{exo:b3:fouriertransform:8}
Partiële integratie ($f \in \mathcal S$; de randtermen
verdwijnen):
\[
1 = \int f^2 = \bigl[xf^2\bigr]_{-\infty}^{\infty} - \int
x\,(f^2)' = -2\int xff' \leq 2\,\norm{xf}_2\,\norm{f'}_2 .
\]
Plancherel en $\widehat{f'} = \iu\xi\hat f$ geven
$\norm{f'}_2^2 = \frac1{2\pi}\int\xi^2\abs{\hat f}^2$.
Kwadrateren van het bovenstaande:
\[
\frac14 \leq \norm{xf}_2^2\cdot\frac1{2\pi}
\int\xi^2\abs{\hat f}^2\dd\xi .
\]
Gelijkheid vereist gelijkheid in Cauchy--Schwarz: $f' = \lambda
xf$ met $\lambda < 0$ (\omterm{def:b3:lebesgue:l1}{integreerbaarheid}), dat wil zeggen $f(x) =
c\,\eu^{\lambda x^2/2}$: de gaussische functies. Een signaal dat
in $x$ geconcentreerd is (kleine $\norm{xf}_2$) moet een
uitgesmeerd spectrum hebben, en omgekeerd: het
onzekerheidsbeginsel.
\end{solution}

\begin{solution}{exo:b3:fouriertransform:9}
De gaussische functie $f(x) = \eu^{-\pi tx^2}$ ligt in $\mathcal
S$, dus is \cref{thm:b3:fouriertransform:poisson} van toepassing,
en $\hat f(\xi) = t^{-1/2}\eu^{-\xi^2/4\pi t}$; in $\xi = 2\pi k$
wordt het rechterlid $t^{-1/2}\eu^{-\pi k^2/t}$: de identiteit
van de thètafunctie. Bij $t = 1$:
\[
\sum_{n\in\Z}\eu^{-\pi n^2} = 1 + 2\eu^{-\pi} + 2\eu^{-4\pi} +
\cdots \approx 1 + 0.0864278 + 0.0000070 = 1.0864348,
\]
nauwkeurig tot op $6$ decimalen met drie termen ($\eu^{-9\pi}
\approx 5\cdot10^{-13}$). Bij $t = 10^{-2}$: de definiërende
reeks vereist $\eu^{-\pi n^2/100} < 10^{-7}$, dus $n \gtrsim 23$
--- ongeveer $47$ termen --- terwijl de getransformeerde reeks
$10\sum_k\eu^{-100\pi k^2}$ is, waarin de term voor $k = 1$ al
$\sim 10^{-136}$ is: één term volstaat.
\end{solution}

\begin{solution}{exo:b3:fouriertransform:10}
$\mathcal Ff \in L^2(\intcc{-\pi}\pi) \subseteq
L^1(\intcc{-\pi}\pi)$ (eindige \omterm{def:b3:measure:measure}{maat}), dus geeft de inversie de
\omterm{def:b3:topology:continuity}{continue} vertegenwoordiger $f(x) =
\frac1{2\pi}\int_{-\pi}^\pi\mathcal Ff(\xi)\eu^{\iu
x\xi}\dd\xi$, met
\[
f(n) = \frac1{2\pi}\int_{-\pi}^{\pi}\mathcal
Ff(\xi)\,\eu^{\iu n\xi}\dd\xi = \langle e_{-n}, \mathcal
Ff\rangle
\]
in de notatie van \cref{thm:b3:hilbert:fourier}. Ontwikkeling in
die \omterm{def:b3:hilbert:onb}{hilbertbasis}: $\mathcal Ff = \sum_nf(n)\,\eu^{-\iu n\xi}$ in
$L^2(\intcc{-\pi}\pi)$. Pas de $L^2$-continue $\mathcal F^{-1}$
term voor term toe:
\[
\mathcal F^{-1}\bigl(\mathbf
1_{\intcc{-\pi}\pi}\eu^{-\iu n\xi}\bigr)(x)
= \frac1{2\pi}\int_{-\pi}^{\pi}\eu^{\iu\xi(x - n)}\dd\xi
= \frac{\sin\bigl(\pi(x-n)\bigr)}{\pi(x - n)} ,
\]
wat $f = \sum_nf(n)\operatorname{sinc}(\cdot - n)$ in $L^2$
geeft: een bandbegrensd signaal ligt vast door zijn monsters in
de gehele getallen --- de bemonsteringsstelling van Shannon.
\end{solution}

\begin{solution}{exo:b3:fouriertransform:11}
(a) Inversie op $\mathcal S$: $f(x) = \frac1{2\pi}\int\hat
f(\xi)\eu^{\iu x\xi}\dd\xi = \frac1{2\pi}(\mathcal F\hat f)(-x)$,
dat wil zeggen $(\mathcal F^2f)(x) = 2\pi f(-x)$. Tweemaal
toegepast: $\mathcal F^4f = 2\pi\,\mathcal F^2f(-\cdot) =
(2\pi)^2f$.

(b) Is $\mathcal Ff = \lambda f$ met $f \neq 0$, dan is
$(2\pi)^2f = \mathcal F^4f = \lambda^4f$, dus $\lambda^4 =
(2\pi)^2$: $\lambda \in \{\pm\sqrt{2\pi}, \pm\iu\sqrt{2\pi}\}$.
De gaussische functie $g(x) = \eu^{-x^2/2}$ heeft $\hat g =
\sqrt{2\pi}\,g$ (\cref{ex:b3:fouriertransform:gaussian} bij $a =
\frac12$): een eigenfunctie voor $+\sqrt{2\pi}$.

(c) $\mathcal F^2f = 2\pi f(-\cdot)$ is $\pm2\pi f$ naargelang de
pariteit. Voor $h(x) = x\eu^{-x^2/2}$: differentiëren van $\hat
g(\xi) = \sqrt{2\pi}\eu^{-\xi^2/2}$ met de regel $\widehat{xf} =
\iu\frac{\dd}{\dd\xi}\hat f$ geeft
\[
\hat h(\xi) = \iu\,\frac{\dd}{\dd\xi}\bigl(\sqrt{2\pi}
\eu^{-\xi^2/2}\bigr) = -\iu\sqrt{2\pi}\,\xi\eu^{-\xi^2/2}
= -\iu\sqrt{2\pi}\,h(\xi) :
\]
een eigenfunctie voor $-\iu\sqrt{2\pi}$. (De hermitefuncties
zetten het patroon voort en doorlopen cyclisch de vier
eigenwaarden --- de discrete klok van Fourier.)

\end{solution}

\begin{solution}{exo:b3:fouriertransform:12}
(a) $\tilde f \in L^2$ met $\norm{\tilde f}_2 = \norm f_2$;
\cref{exo:b3:lp:6} (toegevoegde exponenten $p = q = 2$) maakt
$A_f = f * \tilde f$ begrensd en uniform \omterm{def:b3:topology:continuity}{continu}, met
\[
A_f(x) = \int f(y)\,\overline{f(y - x)}\,\dd y,
\qquad A_f(0) = \norm f_2^2,
\qquad \abs{A_f(x)} \leq \norm f_2\,\norm{f(\cdot -
x)}_2 = A_f(0)
\]
volgens Cauchy--Schwarz.

(b) Voor $f \in L^1\cap L^2$ ligt ook $\tilde f$ in $L^1$, en de
convolutiestelling geeft $\widehat{A_f} = \hat f\,
\widehat{\tilde f}$; en de berekening $\widehat{\tilde f}(\xi) =
\int\overline{f(-x)}\eu^{-\iu\xi x}\dd x = \overline{\int
f(u)\eu^{-\iu\xi u}\dd u} = \overline{\hat f(\xi)}$ geeft
$\widehat{A_f} = \abs{\hat f}^2 \geq 0$.

(c) Is $\abs{\hat f}^2 \in L^1$, dan is de inversie van
toepassing op de \omterm{def:b3:topology:continuity}{continue} $A_f$ (haar getransformeerde is
\omterm{def:b3:lebesgue:l1}{integreerbaar}; \cref{thm:b3:fouriertransform:inversion}):
\[
A_f(x) = \frac1{2\pi}\int\abs{\hat f(\xi)}^2
\eu^{\iu x\xi}\,\dd\xi .
\]
Voor $f = \mathbf 1_{\intcc{-1/2}{1/2}}$ is $\hat f(\xi) =
\frac{2\sin(\xi/2)}\xi = \frac{\sin(\xi/2)}{\xi/2}$, en in $x =
0$:
\[
1 = \norm f_2^2 = \frac1{2\pi}\int_\R
\Bigl(\frac{\sin(\xi/2)}{\xi/2}\Bigr)^2\dd\xi
= \frac1{2\pi}\cdot2\int_\R\Bigl(\frac{\sin u}u\Bigr)^2\dd u
\]
($\xi = 2u$), dat wil zeggen $\int_\R\bigl(\frac{\sin
u}u\bigr)^2\dd u = \pi$ --- de lievelingsintegraal van
Plancherel, teruggevonden met de autocorrelatie.
\end{solution}

\begin{solution}{pb:b3:fouriertransform:1}
\textbf{1.} De vergelijking in $x$ transformeren geeft (formeel)
$\partial_t\hat u(t,\xi) = \widehat{\partial^2_{xx}u} =
(\iu\xi)^2\hat u = -\xi^2\hat u$, voor elke frequentie een
differentiaalvergelijking in $t$: $\hat u(t,\xi) =
\eu^{-t\xi^2}\hat f(\xi)$. Omdat $\eu^{-t\xi^2} = \hat g_t$
(\cref{exo:b3:fouriertransform:4}), is het product de
getransformeerde van $g_t * f$.

\textbf{2.} $\abs{u(t,x)} \leq \norm f_\infty\int g_t = \norm
f_\infty$: goed gedefinieerd. Op $[t_0, T]\times[-A, A]$ is elke
gemengde afgeleide $\partial^m_t\partial^n_xg_t(x - y)$ een
veelterm in $(x - y)$ en $t^{-1}$ maal $\eu^{-(x-y)^2/4t}$, voor
$\abs y \geq 2A$ begrensd door $C\,(1 + y^2)^N\eu^{-(\abs y -
A)^2/4T}$, een \omterm{def:b3:lebesgue:l1}{integreerbare} dominant die op het venster niet van
$(t, x)$ afhangt (en begrensd voor $\abs y \leq 2A$): herhaald
differentiëren onder de integraal
(\cref{thm:b3:lebesgue:paramdiff}) is dus van toepassing, en $u
\in \mathcal C^\infty(\intoo0\infty\times\R)$.

\textbf{3.} Met $g_t(x) = \frac1{2\sqrt{\pi t}}\eu^{-x^2/4t}$:
\[
\partial_tg_t = g_t\Bigl(\frac{x^2}{4t^2} - \frac1{2t}\Bigr)
= \partial^2_{xx}g_t
\]
(tweemaal differentiëren naar $x$: $\partial_xg_t = -\frac
x{2t}g_t$ en $\partial^2_{xx}g_t = \bigl(\frac{x^2}{4t^2} -
\frac1{2t}\bigr)g_t$). Volgens vraag 2 gaan de afgeleiden onder
de integraal door: $\partial_tu = \partial^2_{xx}u$.

\textbf{4.} $u(t,x) - f(x) = \int g_t(y)\bigl(f(x - y) -
f(x)\bigr)\dd y$. Gegeven een \omterm{def:b3:topology:compact}{compacte} $K$ en $\varepsilon$: de
uniforme \omterm{def:b3:topology:continuity}{continuïteit} van $f$ op een \omterm{def:b3:topology:topology}{omgeving} van $K$ geeft een
$\delta$ met $\abs{f(x-y) - f(x)} < \varepsilon$ voor $x \in K$
en $\abs y \leq \delta$; de staart draagt hoogstens $2\norm
f_\infty\int_{\abs y > \delta}g_t(y)\dd y$ bij, dat wil zeggen
$2\norm f_\infty$ maal de massa voorbij $\delta$, gelijk aan
$\frac2{\sqrt\pi}\int_{\delta/2\sqrt t}^\infty\eu^{-z^2}\dd z \to
0$ als $t \to 0$. Voor $f \in L^p$: $\norm{u(t) - f}_p \leq \int
g_t(y)\norm{\tau_yf - f}_p\dd y$ (Minkowski of Jensen als in
\cref{thm:b3:lp:regularization}), en splits op dezelfde manier
met \cref{thm:b3:lp:density}(3).

\textbf{5.} Het ogenblikkelijke gladstrijken is vraag 2 ($u(t)$
is $\mathcal C^\infty$ voor $t > 0$, zonder dat enige gladheid
van $f$ werd gebruikt). Voor $f = \mathbf 1_{\intoo0\infty}$:
\[
u(t, x) = \int_0^\infty g_t(x - y)\dd y
= \frac1{\sqrt\pi}\int_{-x/2\sqrt t}^{\infty}\eu^{-z^2}\dd z
= \frac12\Bigl(1 +
\operatorname{erf}\Bigl(\frac{x}{2\sqrt t}\Bigr)\Bigr),
\qquad \operatorname{erf}(s) =
\frac2{\sqrt\pi}\int_0^s\eu^{-z^2}\dd z :
\]
een gladgestreken sprong waarvan de overgangszone als $\sqrt t$
breder wordt (profielen bij $t_1 < t_2 < t_3$: steeds vlakkere
hellingen door $(0, \frac12)$).

\textbf{6.} De integrand $g_t(x-y)f(y)$ is $\geq 0$ en de kern is
strikt positief: $u(t,x) = 0$ zou $f = 0$ b.o.\ afdwingen. De
monotonie in $f$ is de monotonie van de integraal. Een plek waar
$f$ op een interval nul is, heeft daar toch $u(t, \cdot) > 0$
voor elke $t > 0$: warmte plant zich met oneindige snelheid voort
(elke positiviteit ergens wordt overal ogenblikkelijk gevoeld).

\textbf{7.} Tonelli ($g_t(x-y)\abs{f(y)}$ is \omterm{def:b3:lebesgue:l1}{integreerbaar} op
$\R^2$): $\int u(t,x)\dd x = \int f(y)\bigl(\int g_t(x - y)\dd
x\bigr)\dd y = \int f$.

\textbf{8.} Plancherel: $2\pi\norm{u(t)}_2^2 =
\int\eu^{-2t\xi^2}\abs{\hat f(\xi)}^2\dd\xi$, niet-stijgend in
$t$ (puntsgewijs), strikt dalend tenzij $\hat f = 0$ b.o.\ (dat
wil zeggen $f = 0$), met limiet $0$ als $t\to\infty$ volgens de
gedomineerde convergentie. En $\norm{u(t)}_\infty \leq
\norm{g_t}_\infty\norm f_1 = \frac{\norm f_1}{2\sqrt{\pi t}} \to
0$.

\textbf{9.} Voor $\varphi \in \mathcal C_c^\infty$ is $t \mapsto
\langle\varphi, \hat u(t)\rangle$ $\mathcal C^1$ met afgeleide
$\langle\varphi, \partial_t\hat u\rangle = \langle\varphi,
\widehat{\partial_{xx}u}\rangle = \langle\xi^2\varphi\dots\rangle$
--- preciezer: $\widehat{\partial^2_{xx}u} = -\xi^2\hat u$ draagt
de vergelijking over. Vervolgens heeft voor b.o.\ $\xi$ de
absoluut \omterm{def:b3:topology:continuity}{continue} functie $t \mapsto \eu^{t\xi^2}\hat u(t,\xi)$
afgeleide $\eu^{t\xi^2}(\xi^2\hat u + \partial_t\hat u) = 0$ in
de geïntegreerde zin: ze is constant, en $t \to 0$ laten gaan
($\hat u(t) \to \hat f$ in $L^2$, b.o.\ langs een deelrij) geeft
$\hat u(t, \xi) = \eu^{-t\xi^2}\hat f(\xi)$ b.o. Twee oplossingen
in de klasse hebben dezelfde getransformeerde: ze zijn gelijk.

\textbf{10.} $f = g_s * h$ met $h \in L^2$ dwingt $\hat f =
\eu^{-s\xi^2}\hat h$ af, dat wil zeggen $\hat h =
\eu^{s\xi^2}\hat f \in L^2$. Neem $\hat f(\xi) = \eu^{-\abs\xi}$:
dan is $f(x) = \frac1\pi\cdot\frac1{1 + x^2}$
(\cref{exo:b3:fouriertransform:1}, inversie), een volmaakt gladde
$L^2$-functie; maar $\eu^{2s\xi^2 - 2\abs\xi} \to \infty$: dus
$\eu^{s\xi^2}\hat f \notin L^2$ voor elke $s > 0$. Het profiel
van Cauchy is \emph{nooit} het resultaat van eerdere diffusie.

\textbf{11.} De warmtehalfgroep vermenigvuldigt de
getransformeerden met $\eu^{-t\xi^2}$, dat nergens verdwijnt:
injectief --- formeel wordt geen enkele informatie vernietigd.
Maar haar beeld bestaat uit functies waarvan de getransformeerden
als $\eu^{-t\xi^2}$ vervallen: een minuscule, \omterm{def:b3:topology:interior}{dichte} maar echte
deelruimte van $L^2$ (vraag 10 toont dat zelfs uitstekende
functies er buiten liggen). Inverteren zou de frequentie $\xi$
met $\eu^{t\xi^2}$ versterken: onbegrensd, en dus onstabiel voor
elke verstoring. Diffusie is onomkeerbaar niet omdat de
afbeelding vergeet, maar omdat haar inverse niet \omterm{def:b3:topology:continuity}{continu} kan zijn
--- een pijl van de tijd, gemaakt van functionaalanalyse.

\textbf{12.} $\hat f \in L^2(\intcc{-\Omega}\Omega) \subseteq
L^1$ (Cauchy--Schwarz op een begrensd interval), dus is $F(x) =
\frac1{2\pi}\int_{-\Omega}^\Omega\hat f(\xi)\eu^{\iu
x\xi}\dd\xi$ overal gedefinieerd, en differentiëren onder de
integraal (gedomineerd door $\Omega^k\abs{\hat f} \in L^1$ op de
band) maakt haar $\mathcal C^\infty$ met overal $\abs{F^{(k)}}
\leq \frac{\Omega^k}{2\pi}\norm{\hat f}_{L^1}$. En $F = f$ b.o.:
beide leden hebben dezelfde getransformeerde, en de transformatie
is injectief op $L^2$
(\cref{thm:b3:fouriertransform:plancherel} en haar
$L^2$-uitbreiding).

\textbf{13.} De exponentiëlen $\xi \mapsto \eu^{-\iu
n\pi\xi/\Omega}$, $n \in \Z$, vormen een \omterm{def:b3:hilbert:onb}{hilbertbasis} van
$L^2(\intcc{-\Omega}\Omega)$ (\cref{thm:b3:hilbert:fourier},
herschaald). De coëfficiënt van $\hat f$ langs de $n$-de is
\[
\frac1{2\Omega}\int_{-\Omega}^\Omega\hat f(\xi)\,
\eu^{\iu n\pi\xi/\Omega}\dd\xi
= \frac{2\pi}{2\Omega}\cdot
\frac1{2\pi}\int_{-\Omega}^{\Omega}\hat f(\xi)\,
\eu^{\iu(n\pi/\Omega)\xi}\dd\xi
= \frac\pi\Omega\,f\Bigl(\frac{n\pi}\Omega\Bigr),
\]
volgens de formule van vraag 12 in $x = \frac{n\pi}\Omega$: de
vermelde ontwikkeling geldt in $L^2$ van de band.

\textbf{14.} Voer de ontwikkeling in in de inversieformule van
vraag 12; het verwisselen van som en integraal is de \omterm{def:b3:topology:continuity}{continuïteit}
van de $L^2$-paring tegen $\frac1{2\pi} \eu^{\iu
x\xi}\mathbf 1_{\abs\xi\leq\Omega}$ (met $L^2$-norm
$\frac{\sqrt{2\Omega}}{2\pi}$, onafhankelijk van $x$ --- vandaar
de uniformiteit):
\[
f(x) = \sum_nf\Bigl(\frac{n\pi}\Omega\Bigr)\cdot
\frac1{2\Omega}\int_{-\Omega}^\Omega
\eu^{\iu(x - n\pi/\Omega)\xi}\dd\xi
= \sum_nf\Bigl(\frac{n\pi}\Omega\Bigr)
\operatorname{sinc}(\Omega x - n\pi),
\]
want $\frac1{2\Omega}\int_{-\Omega}^\Omega\eu^{\iu u\xi}\dd\xi =
\frac{\sin(\Omega u)}{\Omega u}$.

\textbf{15.} De berekening van vraag 14 achterstevoren gelezen
geeft dat de getransformeerde van $s_n =
\operatorname{sinc}(\Omega\cdot - n\pi)$ gelijk is aan $\hat s_n
= \frac\pi\Omega\,\eu^{-\iu n\pi\xi/\Omega}\,\mathbf
1_{\intcc{-\Omega}\Omega}$. Plancherel:
\[
\langle s_n, s_m\rangle = \frac1{2\pi}
\Bigl(\frac\pi\Omega\Bigr)^2\int_{-\Omega}^\Omega
\eu^{\iu(n-m)\pi\xi/\Omega}\dd\xi =
\frac\pi\Omega\,\delta_{nm} :
\]
een orthogonale familie met constante norm
$\sqrt{\pi/\Omega}$. Normen nemen in de ontwikkeling van vraag
14: $\norm f_2^2 = \frac\pi\Omega\sum_n\abs{f(n\pi/\Omega)}^2$.

\textbf{16.} $g(x) = \sin(\Omega x)\operatorname{sinc} (\Omega x)
= \frac{\sin^2(\Omega x)}{\Omega x}$ verdwijnt in elk roosterpunt
$\frac{n\pi}\Omega$ (ook in $0$, met de limiet) en is niet
identiek nul. Haar band: schrijf $g = \frac1{2\iu}\bigl(\eu^{\iu
\Omega x} - \eu^{-\iu\Omega x}\bigr)\operatorname{sinc}(\Omega
x)$; modulatie met $\eu^{\pm\iu\Omega x}$ verschuift de
getransformeerde over $\mp\Omega$, dus heeft $\hat g$ drager in
$\intcc{-2\Omega}{2\Omega}$ (zelfs in de vereniging van twee
verschoven banden): $g \in PW_{2\Omega}$, onzichtbaar voor
bemonstering met snelheid $\Omega$ --- aliasing in levenden
lijve.

\textbf{17.} Volgens vraag 15 dragen de monsters de energie
democratisch: $\norm f^2 = \frac\pi\Omega\sum
\abs{f(n\pi/\Omega)}^2$. Is de energie van het signaal buiten het
tijdvenster $\intcc{-T/2}{T/2}$ hoogstens $\varepsilon^2\norm
f^2$, dan voldoen de monsters buiten het venster (op randtermen
na, beheerst door de uniforme afschatting van vraag 12) aan
$\frac\pi\Omega\sum_{\abs{n\pi/\Omega} > T/2}
\abs{f(n\pi/\Omega)}^2 \approx \norm{f\,\mathbf 1_{\abs x >
T/2}}^2 \leq \varepsilon^2\norm f^2$: de bemonsteringsreeks
afknotten tot de ongeveer $\frac{\Omega T}\pi$ indices binnen het
venster reconstrueert $f$ op een relatieve fout $\approx
\varepsilon$ na. Bijgevolg telt het product van tijd en bandbreedte
$\frac{\Omega T}{\pi}$ de effectieve reële vrijheidsgraden van
het signaal --- de regel achter elk audioformaat.

\textbf{18.} (a) $\operatorname{sinc}(\Omega x)$ heeft monsters
$f(n\pi/\Omega) = \operatorname{sinc}(n\pi) = \delta_{n0}$: de
reeks reduceert tot haar term voor $n = 0$,
$\operatorname{sinc}(\Omega x)$ --- de stelling reproduceert haar
eigen kern. (b) Heeft $\hat f$ drager in
$\intcc{-\Omega'}{\Omega'} \subseteq \intcc{-\Omega}\Omega$, dan
loopt elke stap van de vragen 13 en 14 woordelijk door met de
grotere band $\Omega$ (de ontwikkeling van $\hat f$ op het
grotere interval blijft geoorloofd): sneller bemonsteren dan de
eigen nyquistsnelheid verandert niets aan de reconstructie ---
overbemonstering is onschadelijk, en in de praktijk gunstig (er
kunnen dan sneller vervallende reconstructiekernen worden
gebruikt).

\textbf{19.} Ontwikkel $\eu^{-\iu\xi x} = \sum_k\frac{(-\iu\xi
x)^k}{k!}$ binnen de integraal; op $\intcc{-A}A$ convergeert de
reeks normaal ($\sum_k\frac{\abs{\xi}^kA^k}{k!}\abs f \in L^1$),
dus is term voor term integreren geoorloofd:
\[
\hat f(\xi) = \sum_{k\geq0}\frac{(-\iu\xi)^k}{k!}m_k,
\qquad \abs{m_k} \leq A^k\norm f_1 .
\]
Die afschatting laat de reeks voor elke complexe $\xi$
convergeren; rond elk punt $\xi_0$ geeft hergroeperen (absolute
convergentie) een machtreeks in $\xi - \xi_0$: dus is $\hat f$
overal reëel-analytisch met oneindige straal.

\textbf{20.} Zij $g$ reëel-analytisch op $\R$ (haar taylorreeks
convergeert nabij elk punt naar $g$) en $Z = \{\xi :
g^{(k)}(\xi) = 0\ \forall k\}$. $Z$ is gesloten (doorsnede van
gesloten verzamelingen); ze is open, want in $\xi_0 \in Z$ is de
lokale taylorontwikkeling van $g$ de nulreeks, zodat $g$ nabij
$\xi_0$ identiek verdwijnt, samen met al haar afgeleiden.
Verdwijnt $g$ op een interval, dan is $Z \neq \varnothing$; en
wegens de \omterm{def:b3:topology:connected}{samenhang} van $\R$ is $Z = \R$: dus $g \equiv 0$. Had
nu een $f \neq 0$ \omterm{def:b3:topology:compact}{compacte} drager samen met $\hat f$: vraag 19
maakt $\hat f$ reëel-analytisch, en ze verdwijnt buiten een
\omterm{def:b3:topology:compact}{compacte} verzameling, dus op intervallen: bijgevolg $\hat f
\equiv 0$, zodat $f = 0$ b.o.\ wegens de injectiviteit ---
tegenspraak. Evenzo kan een $f \neq 0$ in $PW_\Omega$ geen
\omterm{def:b3:topology:compact}{compacte} drager hebben (verwissel de rollen van $f$ en $\hat f$
met de inversie): bandbegrensde signalen sterven nooit,
tijdbegrensde signalen bezetten een onbegrensd spectrum.

\textbf{21.} Voor $f = \eu^{-ax^2}$ is $\norm f_2^2 =
\sqrt{\frac\pi{2a}}$ en $\int x^2\abs f^2 =
\frac1{4a}\sqrt{\frac{\pi}{2a}}$ (het tweede moment van de
gaussische functie); verder is $\hat f = \sqrt{\frac\pi
a}\,\eu^{-\xi^2/4a}$ (\cref{ex:b3:fouriertransform:gaussian}) en
\[
\frac1{2\pi}\int\xi^2\abs{\hat f}^2\dd\xi =
\frac1{2\pi}\cdot\frac\pi a\int\xi^2
\eu^{-\xi^2/2a}\dd\xi = \frac1{2a}\cdot a\sqrt{2\pi a}
= \frac{\sqrt{2\pi a}}2 .
\]
Genormeerd product: $\frac1{4a}\sqrt{\frac\pi{2a}}\cdot
\frac{\sqrt{2\pi a}}2\big/\frac{\pi}{2a} = \frac14$, onafhankelijk
van $a$. De schaling verklaart die constantheid: $f$ door
$f(\lambda\cdot)$ vervangen vermenigvuldigt $\int x^2\abs
f^2/\norm f^2$ met $\lambda^{-2}$ en $\frac1{2\pi}\int\xi^2\abs{\hat
f}^2/\norm f^2$ met $\lambda^{2}$: het product is een
dilatatie-invariant, en de gaussische functies vormen één
dilatatiebaan.

\textbf{22.} Met $\abs f^2$ als positiedichtheid en
$\frac1{2\pi}\abs{\hat f}^2$ als impulsdichtheid van een
kwantumtoestand (de fysische eenheden voeren $\hbar$ in) luidt
\cref{exo:b3:fouriertransform:8} $\sigma_x\sigma_p \geq
\frac\hbar2$: geen enkele toestand is scherp in beide
observabelen. Doorheen het hoofdstuk draagt één wet vijf pakken:
warmte strijkt ogenblikkelijk glad omdat $\eu^{-t\xi^2}$ de hoge
frequenties vernietigt (Deel II); de stroming kan niet
achterwaarts lopen omdat ze herstellen onbegrensd is (Deel IV);
een bandbegrensd signaal is stijf genoeg om op een aftelbaar
rooster te leven (Deel V); geen enkele functie duikt onder de
bodem van Heisenberg; en geen enkele functie heeft aan beide
zijden van de transformatie \omterm{def:b3:topology:compact}{compacte} drager (vragen 19 en 20).
Wat $\hat f$ op oneindig doet, regeert wat $f$ waar dan ook mag
doen.

\textbf{23.} Zowel $g_t$ als $g_s$ ligt in $L^1$ met
$\widehat{g_t}(\xi) = \eu^{-t\xi^2}$ (de berekening van vraag 1),
dus geeft de convolutiestelling $\widehat{g_t * g_s} =
\eu^{-t\xi^2}\eu^{-s\xi^2} = \eu^{-(t+s)\xi^2} =
\widehat{g_{t+s}}$; twee $L^1$-functies met dezelfde
getransformeerde vallen b.o.\ samen (injectiviteit, via de
inversiestelling --- hier zijn beide leden \omterm{def:b3:topology:continuity}{continu}, dus vallen ze
overal samen): $g_t * g_s = g_{t+s}$. Bijgevolg is $u(t + s) =
g_{t+s} * f = g_s * (g_t * f) = g_s * u(t)$ (associativiteit van
de convolutie, Tonelli). Log-convexiteit: zij $N(t) =
\norm{u(t)}_2^2 = \frac1{2\pi}\int \eu^{-2t\xi^2}\abs{\hat
f}^2\dd\xi$ (Plancherel, vraag 8). Voor $t = \frac{t_1 + t_2}2$
schrijf je
\[
\eu^{-2t\xi^2}\abs{\hat f}^2
= \Bigl(\eu^{-2t_1\xi^2}\abs{\hat f}^2\Bigr)^{1/2}
\Bigl(\eu^{-2t_2\xi^2}\abs{\hat f}^2\Bigr)^{1/2},
\]
en Cauchy--Schwarz geeft $N\bigl(\frac{t_1+t_2}2\bigr) \leq
\sqrt{N(t_1)\,N(t_2)}$: dus is $\ln N$ convex in het midden, en
omdat ze \omterm{def:b3:topology:continuity}{continu} is (gedomineerde convergentie in $t$), convex;
en dus ook $\ln\norm{u(t)}_2 = \frac12\ln N(t)$. Verval met een
convexe logaritme: de warmtestroming kan geen energie in een
uitbarsting verliezen om daarna stil te vallen.

\textbf{24.} De momenten van de kern: $\int g_t = 1$ (vraag 7 met
$f = g_s$, of rechtstreeks de gaussische integraal), $\int
x\,g_t(x)\dd x = 0$ (oneven integrand), en, met de substitutie $x
= 2\sqrt t\,v$,
\[
\int_\R x^2g_t(x)\,\dd x
= \frac{4t}{\sqrt\pi}\int_\R v^2\eu^{-v^2}\dd v = 2t .
\]
Substitueer $x = z + y$ in de convolutie en merk op dat
$\iint(\abs z + \abs y)^2g_t(z)f(y)\,\dd z\,\dd y < \infty$ (elk
van $\int\abs z^kg_t$ en $\int\abs y^kf$ is eindig voor $k \leq
2$, met $\abs y \leq \frac{1 + y^2}2$); Fubini en Tonelli zijn
dan van toepassing op de momentintegralen hieronder:
\[
\int x\,u(t,x)\dd x = \iint (z + y)\,g_t(z)f(y)\,\dd z\,\dd
y = 0\cdot\!\int\! f + 1\cdot\!\int\! yf(y)\dd y,
\]
wat de eerste bewering is; en
\[
\iint (z+y)^2g_t(z)f(y)\,\dd z\,\dd y
= 2t\int f + 2\cdot0\cdot\!\int\! yf + \int y^2f(y)\dd y ,
\]
de tweede. Gemiddelden tellen op, varianties tellen op, en de
kern draagt gemiddelde $0$ en variantie $2t$ bij: na tijd $t$
heeft de warmte zich over een breedte van de orde $\sqrt{2t}$
verspreid --- de afstand groeit als de vierkantswortel van de
tijd, het handschrift van de diffusie (en van de brownse paden
van \cref{ch:b3:probability}).

\textbf{25.} Aan de kant van de transformatie: $\hat f(\xi) =
\sqrt\pi\,\eu^{-\xi^2/4}$, dus $\hat u(t,\xi) =
\sqrt\pi\,\eu^{-(t + \frac14)\xi^2}$, en dat is de
getransformeerde van $(1 + 4t)^{-1/2}\exp\bigl(-x^2/(1+4t)\bigr)$
(het gaussische woordenboek $\eu^{-ax^2} \mapsto
\sqrt{\pi/a}\,\eu^{-\xi^2/4a}$ met $a = \frac1{1+4t}$): de
gesloten vorm. Rechtstreekse controle, met $\sigma = 1 + 4t$:
\[
\partial_tu = \sigma^{-1/2}\eu^{-x^2/\sigma}
\Bigl(-\frac2\sigma + \frac{4x^2}{\sigma^2}\Bigr)
= \partial^2_{xx}u ,
\]
beide leden berekend uit $\partial_xu = -\frac{2x}\sigma\,u$.
Behoud: $\int u(t) = \sigma^{-1/2}\sqrt{\pi\sigma} = \sqrt\pi$
voor alle $t$. Dissipatie:
\[
\norm{u(t)}_2^2 = \frac1\sigma\int\eu^{-2x^2/\sigma}\dd x
= \frac1\sigma\sqrt{\frac{\pi\sigma}2}
= \sqrt{\frac\pi2}\,(1+4t)^{-1/2},
\]
dus $\norm{u(t)}_2 = (\pi/2)^{1/4}(1+4t)^{-1/4}$, niet-stijgend
met convexe logaritme (vraag 23); de $L^2$-norm vervalt als
$t^{-1/4}$, precies de helft van de exponent $t^{-1/2}$ van
$\norm{u(t)}_\infty$ --- in overeenstemming met $\norm u_2^2 \leq
\norm u_\infty\norm u_1$ en het behoud van $\norm u_1$.
Variantie: $\int x^2u(t) = \frac1\sigma\cdot
\frac{\sigma^{3/2}\sqrt\pi}2 = \frac{\sqrt\pi}2\,(1 + 4t) = \int
x^2f + 2t\sqrt\pi$, zoals vraag 24 voorspelt ($\int x^2f =
\frac{\sqrt\pi}2$, $\int f = \sqrt\pi$). Bij $t = 6$: $\sigma =
25$, piekhoogte $u(6, 0) = \frac15$ tegenover $u(0,0) = 1$, een
breedteschaal $\sqrt\sigma = 5$ maal de aanvankelijke, en de hele
tijd $\int u = \sqrt\pi \approx 1.7725$: de vlek is vijf keer
lager, vijf keer breder, en er ontbreekt geen enkele calorie.
\end{solution}
